\section{Introduction}
Modern industrial infrastructures rely heavily on rotating machinery, such as wind turbines, aero-engines, and high-speed railways, where components like bearings and gearboxes operate under demanding conditions \cite{zhai2021multi}. The health status of these machines is critical to operational safety and economic efficiency, as unforeseen failures often lead to catastrophic accidents and massive financial losses \cite{lei2020applications}. Consequently, Intelli-gent Fault Diagnosis (IFD) has emerged as a cornerstone of smart manufacturing, leveraging sensor-based condition monitoring to detect anomalies before they escalate. However, as industrial systems grow in complexity, the demand for more robust, autonomous, and high-precision diagnostic frameworks becomes increasingly urgent to ensure the continuity of critical production lines \cite{li2020deep}.

Despite the success of deep learning in fault diagnosis, most existing methods operate under the closed-set assumption, requiring intensive manual labeling and a balanced distribution of training samples across all health states \cite{zhang2022digital}. In real-world industrial scenarios, machinery typically operates in normal states for the vast majority of its lifespan, making fault samples—especially those representing specific early-stage degradation—extremely rare and difficult to collect \cite{liu2018zero}. This creates a significant "long-tail" distribution problem where traditional supervised models fail to generalize to unseen fault categories. To address this data scarcity and the impracticality of exhaustive labeling, the research focus has shifted from standard supervised learning toward Zero-Shot Learning (ZSL), which aims to recognize novel fault classes without any physical training samples by leveraging auxiliary information \cite{wang2019zero, socher2013zero}.

Current ZSL approaches in fault diagnosis primarily rely on semantic embeddings, such as attributes or word vectors (e.g., Word2Vec or GloVe), to bridge the gap between vibration signals and class labels \cite{lampert2014attribute}. However, these methods face three critical limitations. First, manual attribute definition is labor-intensive and highly subjective, requiring domain experts to enumerate physical characteristics for every possible fault type \cite{fu2018zero}. Second, generic word embeddings are trained on web-scale natural language corpora (e.g., Wikipedia), which lack the technical specificity required for industrial diagnostics; for instance, the term "outer race fault" carries a physical semantic meaning in engineering that generic vectors fail to capture deeply \cite{mikolov2013efficient}. Third, existing models often neglect the complex temporal dependencies inherent in high-frequency vibration data, leading to a "domain gap" where the alignment between the visual-like signal features and the abstract semantic space remains brittle \cite{xian2018zero, buchert2020zero}.

To overcome these challenges, this paper proposes the Cross-Modality Semantic Alignment Transformer (CMSA-Trans), a novel zero-shot fault diagnosis framework that leverages Large Language Models (LLMs) to bridge the semantic divide. Unlike traditional methods, CMSA-Trans utilizes the descriptive power of LLMs to generate high-fidelity, domain-specific fault semantics, transforming abstract class labels into detailed technical descriptions. By treating fault diagnosis as a cross-modal alignment task, our framework employs a Transformer-based temporal encoder to extract discriminative features from raw vibration signals, which are then projected into a shared latent space with the LLM-generated semantics. This alignment ensures that the model learns the underlying physical relationships between signal patterns and their linguistic definitions, enabling the recognition of unseen faults through their semantic proximity to known physical phenomena \cite{vaswani2017attention, radford2021learning}.

The primary contributions of this research are four-fold:
\begin{enumerate}
    \item We introduce the CMSA-Trans framework, which reformulates zero-shot fault diagnosis as a cross-modal semantic alignment task, effectively bypassing the need for labeled samples of target fault categories.
    \item We propose a systematic method to utilize LLMs for generating expert-level fault semantics, capturing the nuanced physical characteristics of rotating machinery faults more effectively than traditional word embeddings or manual attributes.
    \item A Transformer-based temporal feature extractor is designed to capture long-range dependencies in vibration signals, providing a more robust representation for alignment with high-dimensional semantic descriptors.
    \item Extensive experiments on multiple open-source and industrial datasets demonstrate that CMSA-Trans significantly outperforms state-of-the-art ZSL methods, achieving high accuracy on unseen fault classes and showing remarkable generalization across different operating conditions.
\end{enumerate}

The remainder of this paper is organized as follows. Section II provides a comprehensive review of the related literature on intelligent fault diagnosis and zero-shot learning. Section III details the proposed CMSA-Trans architecture, including the LLM semantic generation and the cross-modal alignment mechanism. Section IV presents the experimental setup, results, and comparative analysis. Finally, Section V concludes the paper and discusses future research directions.

Furthermore, the integration of Large Language Models (LLMs), such as GPT-4, provides a paradigm shift in semantic feature engineering for industrial applications \cite{bubeck2023sparks}. Unlike static word embeddings that aggregate broad linguistic associations, LLMs can be prompted to synthesize multi-faceted technical descriptions, capturing the underlying physics of a fault (e.g., impact repetition frequencies, localized high-frequency resonances, and specific degradation patterns in time-frequency domains) \cite{brown2020language}. This allows for a richer and more discriminative semantic space, where the distance between fault classes reflects their actual physical mechanisms rather than their linguistic co-occurrence in general-purpose text. By harnessing these LLM-generated semantics, CMSA-Trans provides a data-efficient and highly descriptive alternative to traditional ZSL approaches.

The critical importance of robust cross-modal alignment cannot be overstated. In traditional deep learning models, features extracted from sensor data are typically mapped directly to one-hot encoded labels, which contain no inherent physical information and thus cannot be generalized to unseen fault types \cite{goodfellow2016deep, lecun2015deep}. In contrast, our semantic alignment approach utilizes a projection layer to map the Transformer's high-level signal representations into the same dimensional space as the LLM-derived technical descriptors. This ensures that the model learns to associate specific signal patterns, such as periodic transients or broadband noise increases, with their corresponding physical descriptions. When a new fault type with a unique technical signature occurs, the model can identify it by matching its signal pattern to its linguistic definition, even if it has never encountered a labeled example during training.

Recent advancements in attention-based architectures have revolutionized how we process sequential data in various fields, from natural language processing to computer vision \cite{devlin2018bert, dosovitskiy2020image}. In the context of rotating machinery fault diagnosis, vibration signals are inherently time-variant and exhibit long-range dependencies, where early-stage fault signatures may be obscured by noise or periodic components from normal operation \cite{zhang2021attention}. While Recurrent Neural Networks (RNNs) and Long Short-Term Memory (LSTM) units have been used to handle such data, they often struggle with gradient vanishing or capturing dependencies over very long windows of high-frequency samples \cite{hochreiter1997long}. The Transformer’s self-attention mechanism, however, provides a more effective way to weight different time steps according to their relevance to the diagnostic task, allowing the CMSA-Trans model to focus on the most informative transients and spectral characteristics across the entire signal window.

This shift toward utilizing large-scale pre-trained models for industrial semantic understanding is motivated by the "world knowledge" packaged within these models \cite{zhao2023brain}. While an engineer might define a bearing fault by its physical dimensions, an LLM can provide a comprehensive textual profile that links those dimensions to expected sensor responses, characteristic frequencies, and interaction effects between components. By aligning these rich descriptions with raw data, our approach moves closer to an "expert-informed" diagnostic system that mimics the reasoning of an experienced technician. This is particularly vital in critical infrastructure where failures are rare but catastrophic, and the ability to identify a new failure mode correctly on its first occurrence is essential for preventing downtime and maintaining structural integrity in modern industrial complexes.

Expanding the scope of zero-shot learning to include cross-modality also addresses the "hubness problem" often encountered in semantic embedding spaces, where certain vectors become nearest neighbors to many unrelated items in high-dimensional space \cite{radovanovic2010hubs}. By utilizing a more descriptive and technically grounded semantic space provided by LLMs, the separation between distinct fault modes—such as inner-race versus outer-race bearing defects—becomes much more pronounced. This enhanced class separation in the semantic domain translates directly to improved diagnostic reliability and higher confidence scores for unseen categories. Our validation on a variety of machinery types, including complex planetary gearsets and varying load-speed conditions, further confirms that the CMSA-Trans framework provides a scalable and robust solution for the next generation of intelligent PHM (Prognostics and Health Management) systems.
